\chapter{Дискретная задача оптимального управления}\label{ch:theorem}

% TODO (M4): ~7 страниц. ГЛАВА С ОСНОВНЫМ ТЕОРЕТИЧЕСКИМ РЕЗУЛЬТАТОМ.
% Структура:
%   4.1 Дискретизация waveform: u[k] ∈ U = {VSH1, VSH2, VSL, VCOM, 0}, k = 0..K-1.
%       Связь с регистром 0x32 SSD1680 (153 байта формата Section 6.7).
%   4.2 Функционал стоимости
%       J(u) = Σ_k (α·E[k](u) + β·G(z[k](u)) + γ·TP[k]),
%       где E из формулы Lin 2024 (P_peak = ½·C·ΔV²·f·V_source),
%       G — residual reflectance, TP — длительность фазы.
%   4.3 Hard-constraint: charge balance Σ_k V[k]·TP[k] = 0.
%       Обоснование: предотвращение DC-смещения и износа панели.
%   4.4 Дискретный принцип максимума (Болтянский → Phogat 2017)
%       — гамильтониан H = c(z,u) + λ·F(z,u) + ν·V·TP;
%       — необходимые условия оптимальности на каждом шаге k.
%   4.5 Основная теорема (со ссылкой на нашу нумерацию):
%       Теорема 4.1 (структура оптимума в классе charge-balanced waveforms).
%       Пусть J — строго возрастающая скаляризация (E, G, τ) в неотрицательном
%       ортанте. Тогда оптимум задачи (4.1)-(4.3) в классе charge-balanced
%       waveforms имеет bang-bang структуру с числом активных фаз
%       K_eff ≤ ... , где K_eff зависит только от размерности подпространства
%       допустимых переходов состояний.
%       Доказательство через игольчатые вариации (Phogat 2017, §3).
%   4.6 Следствия:
%       — конструктивная оценка верхней границы K_eff для нашей задачи;
%       — связь с экспериментально наблюдаемой структурой waveforms.

\section*{Заглушка}

Раздел будет наполнен на этапе~M4 в соответствии с Chapter Plan
(см. \texttt{.ai-factory/RESEARCH.md}).
Эскиз доказательства уже подготовлен в \texttt{docs/notes/proof\_sketches.md}
(задача~D4 в плане~M1).
